% !TeX root = 01_intro_ml_polytech.tex
% !TeX encoding = UTF-8
% !TeX spellcheck = en_GB

\begin{frame}{Python Projects Primer}{\includegraphics[height=12pt]{01_slides/img/tool.pdf}~Settings things up locally\dots}
    \begin{center}
        \large
        \textbf{Q:}~How many of you prefer working on Windows in these sessions? How many on Linux?
    \end{center}

    \pause \vfill

    \begin{columns}[T]
        \begin{column}{0.5\textwidth}
            \begin{center}
                \large \textbf{LINUX}

                Open up a \emph{terminal} window:
                \begin{center}
                    \includegraphics[height=24pt]{01_slides/img/gnome_terminal.pdf}
                \end{center}
            \end{center}
        \end{column}
        \hfil
        \begin{column}{0.5\textwidth}
            \begin{center}
                \large \textbf{WINDOWS}

                Search and click on \emph{GIT (Bash)}:
                \begin{center}
                    \includegraphics[height=24pt]{01_slides/img/git_windows.png}
                \end{center}
            \end{center}
        \end{column}
    \end{columns}
\end{frame}

\begin{frame}[fragile]{Python Projects Primer}{\includegraphics[height=12pt]{01_slides/img/tool.pdf}~Practicing some commands\dots}
    \begin{codewithcard}[language=sh]{Terminal}
# This command prints the current "working directory"
pwd
# This command "lists" the files in the directory
ls
# This command "creates" a new directory
mkdir MyNewDirectory
# This command "changes the directory" = moves to a new location
cd MyNewDirectory
# This command creates an empty file
touch AnEmptyFile.txt
# This command shows the content of files
cat AnEmptyFile.txt
# This command removes a file
rm AnEmptyFile.txt
    \end{codewithcard}
\end{frame}

\begin{frame}{Python Projects Primer}{Introducing the practical tools\dots}
    We choose the \href{https://www.python.org/}{\textbf{Python}~\includegraphics[height=12pt]{01_slides/img/python.pdf}} programming language:
    \begin{itemize}
        \item For its \textbf{simplicity} and \textbf{portability};
        \item For its huge \textbf{ecosystem};
        \item Because it is the \textbf{most widely used language} for AI
    \end{itemize}

    \vfill \pause

    What is a Python \textbf{project} made of and what shall we use?
    \begin{itemize}
        \item \textbf{libraries:}
              \begin{itemize}
                  \item \emph{third parties:} code written by others, available to be \textbf{imported} in our code
                  \item \emph{local:} tools used by the main code in our implementation
              \end{itemize}
        \item \textbf{scripts:} pipelines/sequences of operations to perform
        \item \textbf{notebooks:} like scripts, but divided into standalone cells
        \item \textbf{documentation:} logbooks, instructions, experiment notes, etc.
    \end{itemize}
\end{frame}

\begin{frame}[fragile]{Python Projects Primer}{\includegraphics[height=12pt]{01_slides/img/tool.pdf}~Installing a package manager\dots}
    \vfill

    We shall install \href{https://docs.astral.sh/uv/}{\texttt{uv}} to manage all this for us:
    \begin{codewithcard}[language=sh]{Terminal}
# Copy and execute the following line in the previous location
curl -LsSf https://astral.sh/uv/install.sh | sh
    \end{codewithcard}

    \pause
    You must \textbf{close and restart} your terminal, then \textbf{go back} to the previous location:
    \begin{codewithcard}[language=sh]{Terminal}
# Try executing uv once to check that it has been installed
uv
    \end{codewithcard}

    \vfill
\end{frame}

\begin{frame}[fragile]{Python Projects Primer}{\includegraphics[height=12pt]{01_slides/img/tool.pdf}~Creating a first project\dots}
    \vfill

    You can now create your \textbf{first project}:
    \begin{codewithcard}[language=sh]{Terminal}
# Set name and email once to initialise the versioning system
# This is only necessary ONCE, if you have not already
git config --global user.name "Name Surname"
git config --global user.email name.surname@whatever.com

# Create your first project
uv init --name="my-project-name" \
--description="This is where you describe the purpose of the project" \
--build-backend="hatch" \
--python 3.13
    \end{codewithcard}

    \vfill

    Using \textbf{any file browser} that you know
\end{frame}

\begin{frame}{Python Projects Primer}{A quick look at a project\dots}
    \centering
    \begin{tikzpicture}[scale=0.6, every node/.style={scale=0.6}]
    % Source files
    \node[label=right:src/] (src) at (0, 0) {\includesvg[height=12pt]{01_slides/img/folder.svg}};

    % Main module
    \node[label=right:mypackage/, below right=0.25em of src] (main) {\includesvg[width=12pt]{01_slides/img/folder.svg}};
    \node[label=right:\_\_init\_\_.py, below right=0.25em of main] (init) {\includesvg[height=12pt]{01_slides/img/python.svg}};

    % Subpackage
    \node[label=right:mysubpackage1/, below=0.25em of init] (sub1) {\includesvg[height=12pt]{01_slides/img/folder.svg}};
    \node[label=right:\_\_init\_\_.py, below right=0.25em of sub1] (init1) {\includesvg[height=12pt]{01_slides/img/python.svg}};
    \node[label=right:mysubmodule1.py, below=0.25em of init1] (mod11) {\includesvg[height=12pt]{01_slides/img/python.svg}};
    \node[label=right:mysubmodule2.py, below=0.25em of mod11] (mod12) {\includesvg[height=12pt]{01_slides/img/python.svg}};

    \node[label=right:mysubpackage2/, below left=0.25em of mod12] (sub2) {\includesvg[height=12pt]{01_slides/img/folder.svg}};
    \node[label=right:\_\_init\_\_.py, below right=0.25em of sub2] (init2) {\includesvg[height=12pt]{01_slides/img/python.svg}};
    \node[label=right:mysubmodule1.py, below=0.25em of init2] (mod21) {\includesvg[height=12pt]{01_slides/img/python.svg}};

    % Scripts
    \node[label=right:myscript.py, below left=0.25em of mod21] (script) {\includesvg[height=12pt]{01_slides/img/python.svg}};

    % Configurations
    \node[label=right:pyproject.toml, below=0.25em of script, xshift=-4.5em] (toml) {\includesvg[height=12pt]{01_slides/img/toml.svg}};
    \node[label=right:README.md, below=0.25em of toml] (readme) {\includesvg[height=12pt]{01_slides/img/readme.svg}};

    % Explanations
    \node<2->[CEAred, text width=12em, align=center] (init_exp) at (12, -1) {Must always be present! \\ {\tiny (may contain instructions for packaging)}};
    \draw<2->[->, thick, CEAred] (init.east)+(2.5, 0) -- (init_exp.west);
    \draw<2->[->, thick, CEAred] (init1.east)+(2.5, 0) -- (init_exp.west);
    \draw<2->[->, thick, CEAred] (init2.east)+(2.5, 0) -- (init_exp.south west);

    \begin{scope}[on background layer]
        \node<3->[set, fit=(main)(init)(sub1)(init1)(mod11)(mod12)(sub2)(init2)(mod21)(script), minimum height=8cm, minimum width=8cm, xshift=2cm] (mainmod) {};
        \node<4->[set green, fit=(sub1)(init1)(mod11)(mod12), minimum height=3cm, minimum width=6cm, xshift=2cm] (mainsub1) {};
        \node<4->[set green, fit=(sub2)(init2)(mod21), minimum height=2cm, minimum width=6cm, xshift=2cm] (mainsub2) {};
    \end{scope}

    \node<3->[CEAblue, text width=15em, align=center] (mainmod_exp) at (14, -3) {Contains the local library (the code) we develop};
    \draw<3->[->, thick, CEAblue] (mainmod.east) -- (mainmod_exp.west);

    \node<4->[CEAgreen, text width=15em, align=center] (mainsub_exp) at (14, -5) {Contains (semantically) similar code that needs to stay together};
    \draw<4->[->, thick, CEAgreen] (mainsub1.east) -- (mainsub_exp.west);
    \draw<4->[->, thick, CEAgreen] (mainsub2.east) -- (mainsub_exp.west);

    \node<5->[CEAgrey, text width=15em, align=center] (toml_exp) at (15, -7) {Contains information (configurations) for use and deployment};
    \draw<5->[->, thick, CEAgrey] (toml.east)+(2.7, 0) -- (toml_exp.west);

    \node<6->[CEAgrey, text width=15em, align=center] (readme_exp) at (15, -9) {Contains the documentation};
    \draw<6->[->, thick, CEAgrey] (readme.east)+(2.5, 0) -- (readme_exp.west);

\end{tikzpicture}

\end{frame}


\begin{frame}[fragile]{Python Projects Primer}{\includegraphics[height=12pt]{01_slides/img/tool.pdf}~Setting things up\dots}
    \textbf{Move your first steps} -- Prepare your working space:
    \begin{enumerate}
        \item Open the file \texttt{pyproject.toml} and remove the \texttt{[project.scripts]} section
        \item Open the file \texttt{\_\_init\_\_.py} inside the \texttt{src} directory and remove the code
        \item Create a new directory \texttt{notebooks} on the same level as the \texttt{src} directory
        \item Create a new directory \texttt{data} on the same level as the \texttt{src} directory
        \item Open the file \texttt{README.md} file and keep it open somewhere: use it as a log
    \end{enumerate}

    \pause \vfill
    \textbf{Install the dependencies} -- Ready your project:
    \begin{codewithcard}[language=sh]{Terminal}
# Install third party libraries
# This is only required once per project
uv add numpy scipy scikit-learn xgboost matplotlib spectral statsmodels jupyter
# Run the notebook interface
uv run jupyter lab
    \end{codewithcard}
\end{frame}
